\section{A Free Threshold Correction, Worth Most Where the Model Is Worst}
\label{sec:calibration}

In every experiment in this paper, the model's ranking quality has run above its achieved CSI. That
gap is calibration loss. The model orders cells better than any single global cut-off can exploit.
A cell either goes above the threshold or it does not, and if the threshold is in the wrong place
then good ordering is wasted.

The fix costs nothing. For each scene, pick the threshold that makes the \emph{predicted} wet
fraction match the mean observed wet fraction of the \emph{training} scenes. That prior is a number
available without ever looking at the test label, so this is a free correction and not a peek at the
answer. Concretely, if the training scenes are 3\,\% wet on average, take the top 3\,\% of predicted
cells for this scene, whatever the raw probabilities happen to be.

\subsection{The result, and why the two holdout schemes disagree}
Table~\ref{tab:calibration} reports it. On the leave-one-region-out split, which is the deployment
question, calibration is worth $+0.0272$ CSI at 15 areas. That is 3.4 seed standard deviations, and
it captures 73\,\% of the gap between the deployed threshold and an oracle threshold chosen with
knowledge of the test scene. It replicates at 14 areas at $+0.0261$, which is 5.7 seed standard
deviations.

On the per-tile split the same correction \emph{hurts}, at $-0.0155$.

\begin{table}[!t]
\centering
\caption{Threshold calibration, terrain only, 3 seeds each. The oracle-threshold column is the CSI
achievable if the best single threshold were chosen with knowledge of the test scene, so it is an
upper bound rather than a method. The last column is how much of that gap the free correction
recovers.}
\label{tab:calibration}
\begin{tabular}{lcccc}
\toprule
held out & global thr. & calibrated & oracle thr. & gap closed \\
\midrule
region, 15 areas & 0.1474 & \textbf{0.1746} & 0.1845 & \textbf{73\,\%} \\
region, 14 areas & 0.1305 & \textbf{0.1566} & 0.1615 & \textbf{84\,\%} \\
tile, 14 areas   & 0.2698 & 0.2542 & 0.3030 & $-$47\,\% \\
\bottomrule
\end{tabular}
\end{table}

That sign flip is the mechanism, not noise. Hold out one tile and five sibling tiles of the same
city stay in training, so the training wet-fraction prior already describes the held-out tile well
and the swept global threshold is already close to right. Forcing the wet fraction to a prior can
then only move it away. Hold out an entire region and the global threshold is genuinely
miscalibrated for ground the model has never seen, so matching the prior recovers most of what the
ranking was already capable of.

This supersedes an earlier reading of ours. We previously measured $+0.0154$ from this correction
on a three-area per-tile split, called it about one seed standard deviation, and left it as a
suggestion. On the split that actually corresponds to deployment, and with five times the data, it
is a replicated effect.

\subsection{It helps most exactly where the model is weakest}
Table~\ref{tab:calregion} breaks it down by region. Every region gains, and the coastal one gains
most.

\begin{table}[!t]
\centering
\caption{Calibration by held-out region, terrain only, 15 areas, 3 seeds. Bias is the ratio of
predicted wet cells to observed wet cells at the global threshold, so a value below 1 means the
model is under-warning.}
\label{tab:calregion}
\begin{tabular}{lcccc}
\toprule
region held out & global & calibrated & change & bias \\
\midrule
\textbf{Mumbai} (coastal)   & 0.0613 & \textbf{0.1061} & \textbf{+73\,\%} & 0.78 \\
Chennai (coastal)           & 0.1642 & 0.1889 & +15\,\% & 0.81 \\
Patna (inland)              & 0.1537 & 0.1711 & +11\,\% & 0.51 \\
Karnataka (inland)          & 0.2207 & 0.2335 & +5.8\,\% & 0.71 \\
\bottomrule
\end{tabular}
\end{table}

The bias column explains why. At the global threshold the model finds only half to four-fifths of
the water that is actually there. Under-warning is the wrong direction for a flood product, because
the cost of missing a flooded street is not the same as the cost of falsely flagging a dry one.
Ranking degrades far less than CSI does, which is what told us this was calibration rather than a
failure of skill.

Put together with the coverage result of Section~\ref{sec:coastal}, held-out coastal Mumbai has
moved from 0.033, when no coast was in training, to 0.061 once one coast was added, to 0.106 with
the free threshold correction applied. That is 3.2 times the starting figure, and it is roughly
5.7 times the all-wet baseline for that region. Neither step required a larger model or a finer
grid. One was data coverage and the other was arithmetic.
